\chapter{The U.S. freight application}
\label{ch:rsue}

\section{Purpose}

The paper applies the statistic to the U.S. freight network. The reproducible
network has 234 nodes: 228 domestic locations and six foreign regions. The
policy set contains 704 directed Interstate road arcs paired into 352
bidirectional physical links. Each experiment lowers the primitive road cost
by one percent in both directions of one physical link.

The exercise is a sufficient-statistic application. It does not estimate a
new 234-location structural model from raw microdata. The adapter combines
model-ready domestic trade and modal flow matrices, location shares, the road
network, and directional 2017 Census port trade. It then checks and constructs
the objects required by the derivative.

\section{International flows}

The headline vintage preserves separate import and export directions from
public 2017 Census containerized vessel trade. Port values are mapped to
domestic nodes and six foreign regions. Because the maintained theory uses a
balanced-flow system, imports and exports are projected onto common normalized
margins. The projection is an explicit model adapter, not a claim that raw
port-level imports equal exports.

The raw monthly API cache is not committed. The repository includes:
\begin{itemize}
  \item the credential-safe downloader;
  \item port and foreign-region crosswalks;
  \item the aggregate derived overlay;
  \item source metadata and hashes;
  \item diagnostics for coverage and balancing.
\end{itemize}
The API key is read from the environment and is never written to a URL, log,
manifest, or output.

\section{Headline specification}

The accepted configuration uses:
\[
 \alpha=0.10,\quad \beta=-0.30,\quad \sigma=9,\quad
 \eta=1.099,\quad \lambda_{\mathrm{road}}=0.092.
\]
It uses \code{ChoiceLogsum}, all observed transport modes in equilibrium,
edge-local road congestion, and no headline terminal congestion.

The policy is a primitive road-cost improvement. The traditional statistic is
the two-direction road traffic share. The extended statistic includes route,
mode, congestion, and spatial-equilibrium responses.

\section{Accepted aggregate results}

The current package reproduces the accepted paper ledger:
\begin{table}[H]
\centering
\caption{Accepted physical-link result summary}
\label{tab:rsue-summary}
\begin{tabular}{lr}
\toprule
Object & Value \\
\midrule
Physical links & 352 \\
Mean extended elasticity & 0.0002230950 \\
Median extended elasticity & 0.0001807022 \\
Maximum extended elasticity & 0.0008718750 \\
Pearson correlation with traditional statistic & 0.9474377 \\
Spearman rank correlation & 0.9591921 \\
\bottomrule
\end{tabular}
\end{table}

The high correlations show that traffic contains substantial information.
They do not imply identical targeting. Only four links overlap between the
two top-ten lists in the accepted result vintage. The paper reports names only
where endpoints can be assigned by point-in-polygon matching to public 2018
Census CBSA geography.

\section{Result provenance}

The paper artifact builder writes directed and physical results, a claim
ledger, TeX macros, figures, map geometry, sensitivity paths, diagnostics, and
a run manifest. The builder requires:
\begin{itemize}
  \item matching configuration and input hashes;
  \item an accepted nonlinear finite-difference report bound to the derivative
  source hashes;
  \item 234 nodes, 704 directed policies, and 352 reciprocal links;
  \item all algebraic and condition-number gates;
  \item deterministic output hashes.
\end{itemize}

Later package versions add columns to the decomposition tables. The numerical
claim ledger remains unchanged when the economics is unchanged, but byte
hashes of the expanded CSV files must be updated. For this reason the guide
reports both numerical equality and artifact hashes rather than treating one
as a substitute for the other.

\section{Reproducing with external inputs}

Set the external data root and run the locked environment:
\begin{lstlisting}[language=bash]
export RSUE_DATA_ROOT=/path/to/hash-verified/rsue-inputs
julia --project=replication/rsue/environment \
  -e 'using Pkg; Pkg.instantiate()'
julia --project=replication/rsue/environment \
  replication/rsue/verify_choice_logsum_fd.jl \
  replication/rsue/rsue_paper_choice_edge_census_2017.toml \
  replication/rsue/verification/choice_logsum_rsue_fd.toml
julia --project=replication/rsue/environment \
  replication/rsue/build_paper_artifacts.jl
\end{lstlisting}
The manifest rejects files whose hashes differ from the declared data vintage.
There is no fallback to the legacy port matrix.

\section{Interpretation}

The result ranks local model-implied benefits, not project net present values.
A high-ranked link may be expensive to build, face environmental constraints,
or conflict with distributional objectives. Conversely, a low-ranked capacity
improvement may be justified by safety or resilience channels outside the
model. The appropriate use is to add the spatial-equilibrium benefit measure
to a broader appraisal, not to suppress the rest of the appraisal.
